%% toc: P15 170 170.5,170.6
% Repeated fixed-sample tests on accumulating data: dependence between looks and the actual first-crossing error.
\chapter{Repeated Fixed-Sample Tests and Optional-Stopping Failures}
\label{sec:p15-w02b-p15-061-cumulative-looks}

A fixed-sample test is calibrated for one preselected sample size. When the
same threshold is applied to the cumulative statistic after every new batch
of observations, and the analyst rejects at the first look that crosses it,
the procedure actually used is a sequential one. Its type-I error is the
probability that the null path crosses the threshold at some look, not the
level of any single look. For looks at independent samples this probability
is calculated in \cref{thm:p15-se-repeated-fixed}. Here the looks share
data: each cumulative statistic contains every earlier observation. This
chapter states the joint law of the looks, gives the exact first-crossing
error for normal observations with known variance, proves that it exceeds
the nominal level as soon as there are two looks but stays below the error
of the same number of independent looks, and proves that it tends to one
under unbounded inspection for any independent identically distributed
observations with finite positive variance.

\emph{Inputs.} A null model for the observations, a look schedule fixed
before the data are seen, a number of looks (finite or unbounded) and a
threshold. \emph{Conclusions.} The actual type-I error of the rule that
rejects at the first crossing; for a finite number of looks, the constant
threshold that restores a prescribed level.

\section{The procedure and its information}
\label{sec:p15-w02b-p15-061-setting}

\begin{hbdefinition}{def:p15-w02b-p15-061-procedure}{Cumulative looks and the first-crossing rule}
Under the null hypothesis let \(X_1,X_2,\ldots\) be independent and
identically distributed with known mean \(\mu_0\) and known variance
\(\sigma^2\in(0,\infty)\). Put
\[
 Y_i=\frac{X_i-\mu_0}{\sigma},\qquad
 S_n=\sum_{i=1}^nY_i,\qquad
 Z_n=\frac{S_n}{\sqrt n}\quad(n\geq1),
\]
and \(\mathcal F_n=\sigma(X_1,\ldots,X_n)\). A \emph{look schedule}
is a strictly increasing sequence of positive integers
\(n_1<n_2<\cdots\), finite with \(K\) terms or infinite, fixed before
any observation is made. The \emph{one-sided fixed-sample test at sample
size \(n\)} with threshold \(c\in\mathbb R\) rejects when \(Z_n>c\).
The \emph{first-crossing rule} over the schedule inspects
\(Z_{n_1},Z_{n_2},\ldots\) in order and stops at
\[
 T=\min\{k:Z_{n_k}>c\},\qquad \min\emptyset=\infty,
\]
rejecting the null hypothesis if \(T\) is finite (with \(K\) looks: if
\(T\leq K\)). The decision to continue after look \(k\) uses only
\(\mathcal F_{n_k}\), so \(\{T=k\}\in\mathcal F_{n_k}\).
Its \emph{actual type-I errors} are
\[
 \alpha_K(c)=P_0\Bigl(\max_{1\leq k\leq K}Z_{n_k}>c\Bigr)=P_0(T\leq K),
 \qquad
 \alpha_\infty(c)=P_0(T<\infty).
\]
The \emph{two-sided} versions replace \(Z_{n_k}>c\) by
\(|Z_{n_k}|>c\) with \(c>0\); their errors are written
\(\alpha^{\pm}_K(c)\) and \(\alpha^{\pm}_\infty(c)\).
Write \(\Phi\) and \(\varphi\) for the standard normal distribution
function and density and \(\bar\Phi=1-\Phi\).
The \emph{Gaussian null model} is the special case
\(X_i\sim N(\mu_0,\sigma^2)\). In that model a single preselected look
has level exactly \(\bar\Phi(c)\), and \(c=\bar\Phi^{-1}(\alpha)\) gives
the fixed-sample level-\(\alpha\) test.
\end{hbdefinition}

The looks are not independent: \(Z_{n_{k+1}}\) is built from the
observations behind \(Z_{n_k}\) and some new ones. The next result makes
that dependence explicit.

\section{The joint law of the looks}
\label{sec:p15-w02b-p15-061-dependence}

\begin{hblemma}{thm:p15-w02b-p15-061-look-dependence}{proved=annex:p15-w02b-p15-061-look-dependence}{Overlapping looks are correlated and Markov}
Under the hypotheses of \cref{def:p15-w02b-p15-061-procedure}, for every
schedule and all \(j\leq k\),
\[
 \operatorname{Cov}_0(Z_{n_j},Z_{n_k})=\sqrt{n_j/n_k}.
\]
Put \(\rho_k=\sqrt{n_k/n_{k+1}}\in(0,1)\) and
\[
 W_{k+1}=\frac{S_{n_{k+1}}-S_{n_k}}{\sqrt{n_{k+1}-n_k}}.
\]
Then, for every \(k\),
\[
 Z_{n_{k+1}}=\rho_kZ_{n_k}+\sqrt{1-\rho_k^2}\,W_{k+1},
\]
the variables \(Z_{n_1},W_2,W_3,\ldots\) are independent, each with mean
\(0\) and variance \(1\), and \(W_{k+1}\) is independent of
\(\mathcal F_{n_k}\).
In the Gaussian null model these variables are independent standard
normal, so \((Z_{n_1},\ldots,Z_{n_K})\) is centered normal with the
displayed covariances and has the density
\[
 f_K(z_1,\ldots,z_K)=\varphi(z_1)\prod_{k=1}^{K-1}
 \frac1{\sqrt{1-\rho_k^2}}\,
 \varphi\Bigl(\frac{z_{k+1}-\rho_kz_k}{\sqrt{1-\rho_k^2}}\Bigr),
\]
which is strictly positive on \(\mathbb R^K\).
\end{hblemma}

The correlation between two looks depends only on the ratio of their sample
sizes. Looks at \(n,2n,\ldots,Kn\) therefore have the same joint law for
every \(n\).

\section{The actual first-crossing error}
\label{sec:p15-w02b-p15-061-error}

\begin{hbtheorem}{thm:p15-w02b-p15-061-finite-error}{proved=annex:p15-w02b-p15-061-finite-error}{Exact error of a repeated one-sided test with finitely many looks}
Assume the Gaussian null model and a schedule
\(n_1<\cdots<n_K\), and let \(c\in\mathbb R\).
\begin{enumerate}
\item \(\displaystyle 1-\alpha_K(c)=\int_{(-\infty,c]^K}f_K(z)\,dz.\)
\item Equivalently, put \(h_K\equiv1\) and, for \(k=K-1,\ldots,1\),
\[
 h_k(z)=\int_{-\infty}^{c}\frac1{\sqrt{1-\rho_k^2}}\,
 \varphi\Bigl(\frac{u-\rho_kz}{\sqrt{1-\rho_k^2}}\Bigr)h_{k+1}(u)\,du.
\]
Then \(\alpha_K(c)=1-\int_{-\infty}^{c}\varphi(z)h_1(z)\,dz\).
\item With two looks and \(\rho=\rho_1=\sqrt{n_1/n_2}\),
\[
 \alpha_2(c)=\bar\Phi(c)+\int_{-\infty}^{c}\varphi(z)\,
 \bar\Phi\Bigl(\frac{c-\rho z}{\sqrt{1-\rho^2}}\Bigr)dz .
\]
\item \(\alpha_1(c)=\bar\Phi(c)\), and for every \(K\geq2\)
\[
 \alpha_{K-1}(c)<\alpha_K(c)<\min\{1,K\bar\Phi(c)\}.
\]
In particular, with \(c=\bar\Phi^{-1}(\alpha)\) and \(0<\alpha<1\),
every single preselected look has level exactly \(\alpha\), while the
first-crossing rule over two or more looks has type-I error strictly
greater than \(\alpha\).
\item A rule that inspects only some of the looks \(n_1,\ldots,n_K\),
chosen in any way from the data already seen, and rejects at the first
inspected look with \(Z_{n_k}>c\), has type-I error at most
\(\alpha_K(c)\).
\end{enumerate}
\end{hbtheorem}

Overlap makes the looks positively dependent, and this dependence works
against the analyst less than independence would. The comparison rests on
the following inequality for monotone functions of independent variables.

\begin{hblemma}{thm:p15-w02b-p15-061-association}{proved=annex:p15-w02b-p15-061-association}{Monotone functions of independent variables are positively correlated}
Let \(U_1,\ldots,U_m\) be independent real random variables and let
\(f,g\colon\mathbb R^m\to[0,1]\) be Borel functions, each nonincreasing in
every coordinate when the other coordinates are held fixed. Then, with
\(U=(U_1,\ldots,U_m)\),
\[
 E\bigl[f(U)g(U)\bigr]\geq E\bigl[f(U)\bigr]\,E\bigl[g(U)\bigr].
\]
\end{hblemma}

\begin{hbtheorem}{thm:p15-w02b-p15-061-versus-independent}{proved=annex:p15-w02b-p15-061-versus-independent}{Overlapping versus independent looks}
Assume the hypotheses of \cref{def:p15-w02b-p15-061-procedure} and a
schedule \(n_1<\cdots<n_K\). For every \(c\in\mathbb R\),
\[
 1-\alpha_K(c)\geq\prod_{k=1}^KP_0\bigl(Z_{n_k}\leq c\bigr).
\]
In the Gaussian null model with \(K\geq2\) the inequality is strict, so
\[
 \bar\Phi(c)<\alpha_K(c)<1-\bigl(1-\bar\Phi(c)\bigr)^K .
\]
With \(c=\bar\Phi^{-1}(\alpha)\) and \(0<\alpha<1\) this reads
\(\alpha<\alpha_K(c)<1-(1-\alpha)^K\); the right side is the error of
the first-crossing rule over \(K\) independent looks of level \(\alpha\)
in \cref{thm:p15-se-repeated-fixed}.
\end{hbtheorem}

The theorem uses the fact that each one-sided acceptance event
\(\{Z_{n_k}\leq c\}\) shrinks when any block of observations increases.
The two-sided acceptance events \(\{|Z_{n_k}|\leq c\}\) are not monotone
in this sense, and the theorem makes no comparison for the two-sided rule;
\cref{def:p15-w02b-p15-061-example} reports two-sided values computed from
the exact integral.

\begin{hbtheorem}{thm:p15-w02b-p15-061-unbounded}{proved=annex:p15-w02b-p15-061-unbounded}{Unbounded inspection rejects a true null almost surely}
Assume only the hypotheses of \cref{def:p15-w02b-p15-061-procedure}
(independent identically distributed observations with finite positive
variance; normality is not required), and let the schedule be infinite.
Then for every \(c\in\mathbb R\)
\[
 P_0\bigl(Z_{n_k}>c\ \text{for infinitely many}\ k\bigr)=1 .
\]
Consequently \(\alpha_\infty(c)=1\) and \(\alpha_K(c)\uparrow1\) as
\(K\to\infty\), for every threshold \(c\), however small the
fixed-sample level \(\bar\Phi(c)\) it was chosen for.
\end{hbtheorem}

For the schedule \(n_k=k\) the conclusion is part of
\cref{thm:p15-clt-oscillation}; the theorem above holds for every
infinite schedule fixed in advance, however sparse.

\begin{hbcorollary}{thm:p15-w02b-p15-061-two-sided}{proved=annex:p15-w02b-p15-061-two-sided}{The two-sided repeated test}
Let \(c>0\). In the Gaussian null model,
\(\alpha^{\pm}_1(c)=2\bar\Phi(c)\),
\[
 1-\alpha^{\pm}_K(c)=\int_{[-c,c]^K}f_K(z)\,dz,
\]
the recursion of \cref{thm:p15-w02b-p15-061-finite-error} holds with
\((-\infty,c]\) replaced by \([-c,c]\) throughout, and for \(K\geq2\)
\[
 \alpha^{\pm}_{K-1}(c)<\alpha^{\pm}_K(c)<\min\{1,2K\bar\Phi(c)\}.
\]
Under the hypotheses of \cref{thm:p15-w02b-p15-061-unbounded} and an
infinite schedule, \(\alpha^{\pm}_\infty(c)=1\).
\end{hbcorollary}

\section{Calibrating a constant threshold}
\label{sec:p15-w02b-p15-061-calibration}

\begin{hbtheorem}{thm:p15-w02b-p15-061-calibrated}{proved=annex:p15-w02b-p15-061-calibrated}{A constant threshold with exact level exists only for finitely many looks}
Assume the Gaussian null model, a schedule \(n_1<\cdots<n_K\) and
\(0<\alpha<1\). The map \(c\mapsto\alpha_K(c)\) is continuous and
strictly decreasing from \(1\) to \(0\), so there is exactly one
\(c_K\in\mathbb R\) with \(\alpha_K(c_K)=\alpha\). Here
\(c_1=\bar\Phi^{-1}(\alpha)\) and \(c_{K}>c_{K-1}\) for \(K\geq2\),
where \(c_{K-1}\) is computed for the first \(K-1\) looks.
The first-crossing rule with threshold \(c_K\) at the looks
\(n_1,\ldots,n_K\) has type-I error exactly \(\alpha\).
With an infinite schedule no constant threshold has error below one, by
\cref{thm:p15-w02b-p15-061-unbounded}.
\end{hbtheorem}

Over an unbounded horizon a level guarantee therefore needs a boundary
that changes with the look, or a different statistic. For two simple
hypotheses the sequential probability ratio test supplies one: its
upper-exit error over the entire horizon is controlled by the bounds of
\cref{thm:p15-se-errors}.

\begin{hbmethod}{def:p15-w02b-p15-061-method}{Stating the error of a repeated test on accumulating data}
\begin{enumerate}
\item Fix the look schedule \(n_1<\cdots<n_K\) and the number of looks
before collecting data. If the analyst may skip looks on the basis of
the data, use the full list of permitted looks; part~5 of
\cref{thm:p15-w02b-p15-061-finite-error} then bounds the error.
\item Compute \(\rho_k=\sqrt{n_k/n_{k+1}}\).
\item Either report \(\alpha_K(c)\) for the threshold actually used,
computed from the integral or recursion of
\cref{thm:p15-w02b-p15-061-finite-error} (two-sided: from
\cref{thm:p15-w02b-p15-061-two-sided}), or choose \(c_K\) of
\cref{thm:p15-w02b-p15-061-calibrated} for the target level.
\item Report the result as the error of the whole first-crossing
procedure, not as the level of the look at which rejection happened.
\end{enumerate}
Side conditions: the exact values require normal observations with known
variance and a schedule not chosen in response to the data beyond the
skipping allowed in step~1. With no bound on the number of looks no
constant threshold works (\cref{thm:p15-w02b-p15-061-unbounded}).
\end{hbmethod}

\section{A worked calculation}
\label{sec:p15-w02b-p15-061-worked}

\begin{hbexample}{def:p15-w02b-p15-061-example}{Five-percent thresholds on a growing sample}
Take the Gaussian null model, the one-sided threshold
\(c=\bar\Phi^{-1}(.05)\approx1.644854\), and looks at
\(n,2n,\ldots,Kn\), so \(\rho_k=\sqrt{k/(k+1)}\).
With two looks, \(\rho=1/\sqrt2\) and part~3 of
\cref{thm:p15-w02b-p15-061-finite-error} gives
\[
 \alpha_2(c)=.05+\int_{-\infty}^{c}\varphi(z)\,
 \bar\Phi\bigl(\sqrt2\,c-z\bigr)dz\approx.05+.030076=.080076 .
\]
The second term is the probability of not crossing at the first look and
crossing at the second. Evaluating the recursion of part~2 numerically
gives the following first-crossing errors; the right column is the value
\(1-.95^K\) for the same number of independent looks from
\cref{thm:p15-se-repeated-fixed}. By
\cref{thm:p15-w02b-p15-061-versus-independent} the left column lies
strictly between \(.05\) and the right column for every \(K\geq2\).
\begin{center}
\begin{tabular}{rcc}
\(K\) & overlapping looks \(\alpha_K(c)\) & independent looks\\
\hline
1 & .0500 & .0500\\
2 & .0801 & .0975\\
3 & .1011 & .1426\\
4 & .1171 & .1855\\
5 & .1300 & .2262\\
10 & .1718 & .4013\\
20 & .2145 & .6415\\
50 & .2701 & .9231
\end{tabular}
\end{center}
For the two-sided five-percent threshold \(c=\bar\Phi^{-1}(.025)\approx
1.959964\) the same schedule gives
\(\alpha^{\pm}_2\approx.0831\), \(\alpha^{\pm}_5\approx.1417\) and
\(\alpha^{\pm}_{10}\approx.1934\). The one-sided constant thresholds of
\cref{thm:p15-w02b-p15-061-calibrated} that restore level \(.05\) are
\(c_2\approx1.8754\) and \(c_5\approx2.1217\).

The values were obtained by Gauss--Legendre quadrature of the recursion
on \((-12,c]\); restricting every look to \((-12,c]\) changes the
computed probability by at most \(K\Phi(-12)<10^{-31}\) for
\(K\leq50\). For \(K\leq20\) the values agree with an independent
multivariate normal integration to the four decimals shown; at
\(K=50\) refining the quadrature from 1000 to 4000 nodes leaves the
value unchanged to seven decimals. The two-look value was also
computed from the one-dimensional integral in high precision. They are
numerical evaluations of the exact formulas, not simulated frequencies.
The slow growth of the column is consistent with
\cref{thm:p15-w02b-p15-061-unbounded}, which gives the limit one but no
rate; fifty looks still leave the error near \(.27\).
\end{hbexample}

\section{Failure case and limits}
\label{sec:p15-w02b-p15-061-limits}

\begin{hbwarning}{def:p15-w02b-p15-061-warning}{A valid look is not a valid procedure}
Reporting ``significant at the five-percent level'' after rejecting at the
first of several looks that crossed \(\bar\Phi^{-1}(.05)\) misstates the
error of the procedure actually followed: by
\cref{thm:p15-w02b-p15-061-finite-error} its error is strictly larger
whenever two or more looks were possible, and by
\cref{thm:p15-w02b-p15-061-unbounded} an analyst prepared to continue
indefinitely rejects a true null hypothesis with probability one.

The exact values need the Gaussian null model with known variance. With
an estimated variance, a non-normal observation law or a schedule chosen
in response to the data in a way not covered by part~5 of
\cref{thm:p15-w02b-p15-061-finite-error}, the finite-look values above
are not established by these results and require their own calculation.
\Cref{thm:p15-w02b-p15-061-unbounded} needs only independent identically
distributed observations with finite positive variance and a schedule
fixed in advance.
\end{hbwarning}

\emph{Relationships.} \Cref{thm:p15-se-repeated-fixed} is the
independent-look case, and \cref{thm:p15-w02b-p15-061-versus-independent}
places the overlapping-look error below it. \Cref{thm:p15-w02b-p15-061-unbounded} extends the
all-\(n\) oscillation of \cref{thm:p15-clt-oscillation} to arbitrary fixed
schedules, using \cref{thm:p15-kolmogorov-tail-zero-one} and
\cref{thm:p15-clt-iid}. Stopping at a data-dependent time is the setting of
\cref{sec:p15-optional-stopping-and-random-walks}; the general
multiplicity problem of combining several tests is treated in
\cref{sec:p15-multiple-testing}; the sequential probability ratio
construction of \cref{sec:p15-sequential-evidence} gives an
unbounded-horizon error bound that a constant fixed-sample threshold
cannot. When the two hypotheses are replaced by neighborhoods of
distributions, the corresponding sequential test and its expected-size
qualifications at unbounded endpoints are developed in
\cref{sec:p15-robust-hypothesis-testing} and
\cref{sec:p15-robust-sequential-endpoints}.

\emph{Sources.} The normal sums, tail zero-one law and central limit
theorem used are Durrett's Theorem~2.1.20, Example~2.5.2 with
Theorem~2.5.3, and Theorem~3.4.1
\cite[p.~52; p.~81; p.~144]{DurrettPTEAuthor2019}. The dependence
lemma, the inequality for monotone functions of independent variables,
the exact error formulas, the strict inequalities, the comparison with
independent looks, the calibration result and the numerical values are proved or computed in this chapter
and its annex.
\index{sequential testing!repeated fixed-sample test}
\index{optional stopping!cumulative looks}
\index{type-I error!first crossing}
